% sections/04_fault_analysis.tex

\section{Data Quality and Fault Analysis}

The underlying dataset introduces structural biases and limitations that
constrain the conclusions that can responsibly be drawn.
These faults were identified through exploratory visualisation and must be
considered when interpreting every chart in this report.

\begin{table}[H]
\centering\small
\begin{tabular}{@{} c p{4.8cm} p{7.5cm} @{}}
  \toprule
  \textbf{\#} & \textbf{Fault} & \textbf{Primary Impact} \\ \midrule
  1  & Severity Label Imbalance        & Severity 2 $\approx$80\%; misleading raw severity ratios \\
  2  & Geographic Sampling Bias        & CA over-represented; invalid cross-state comparison \\
  3  & Road Network Bias               & Rural areas appear artificially safer \\
  4  & Confounding Variables (Weather) & Weather–severity link is correlational, not causal \\
  5  & Exposure Bias (Time-based)      & No normalisation by vehicles on road per hour \\
  6  & Temporal Data Drift             & Year-on-year trends may reflect data artefacts \\
  7 & Missing Exposure Data           & All severity ratios are accident-conditional \\
  8 & Lack of Contextual Features     & No road type, speed limit, or traffic density \\
  \bottomrule
\end{tabular}
\caption{Summary of identified data faults}
\end{table}

\subsection{Fault 1: Severity Label Imbalance}
Severity~2 dominates (\textasciitilde80--85\% of total data).
Severity~3 and~4 are significantly underrepresented; Severity~1 is extremely
rare. High-severity patterns become statistically weak due to low sample sizes.

\textit{Impact:} All High Severity Rate percentages must be read in the context
of their underlying counts. A small percentage-point shift can be substantial
in absolute terms.

\subsection{Fault 2: Geographic Sampling Bias}
California has disproportionately high accident counts relative to its
population share. Other large states (TX, FL) are significantly lower than
expected. Coverage is dependent on reporting and sensor infrastructure, not
actual accident rates.

\textit{Impact:} Direct state-level comparisons are partially reflections of
data collection coverage rather than true accident rates.

\subsection{Fault 4: Confounding Variables in Weather Analysis}
Weather does not occur in isolation. The analysis lacks simultaneous control for
traffic density, vehicle speed, road type, and time of day.

\textit{Impact:} Causal interpretations of H1 (weather $\rightarrow$ severity)
must be treated as correlational only.

\subsection{Fault 5: Exposure Bias in Time-Based Analysis}
No normalisation by number of vehicles on the road at each hour is applied.
A 21\% severity rate at 4\,am with 1,000 accidents represents a fundamentally
different risk level from 19\% at 8\,am with 50,000 accidents.

\textit{Impact:} Rush-hour severity comparisons (H2) are misleading without
exposure normalisation. All time-based conclusions are accident-conditional.

\subsection{Fault 6: Temporal Data Drift}
Accident counts increase year-on-year; the high-severity rate simultaneously
declines sharply. This decline likely reflects expansion of API coverage
(more Severity~2 records added), not genuine improvements in accident outcomes.

\textit{Impact:} Year-over-year severity rate comparisons must be treated as
exploratory. The 2023 partial year was excluded for this reason.

\subsection{Fault 7: Missing Exposure Data (Critical)}
The dataset contains no information on traffic volume, vehicle miles travelled,
or number of active drivers. True risk per unit of exposure cannot be computed.

\textit{Impact:} The most fundamental limitation for both H1 and H2. All
severity ratios are \emph{accident-conditional}, not \emph{trip-conditional}.

% TODO: Expand remaining fault subsections (3, 6, 8, 9, 11, 12) if required.

\begin{quote}
  \itshape
  While all visualisations accurately represent the data as collected, the
  dataset introduces multiple biases and structural limitations. Any conclusions
  drawn---including the validation of H1 and H2---must be qualified by these
  faults. The analysis demonstrates patterns and associations, not causal
  mechanisms. Deeper investigation using matched datasets (traffic volume, speed
  records, injury severity) would be required to move from correlation to
  causation.
\end{quote}
